\documentclass[../../../uem_phy_std_a.tex]{subfiles}

\begin{document}
    \begin{enumerate}
        \item%
            運ぶ際の電位の変化を$\Delta \phi$とすれば，陽子の位置エネルギー変化%
            $e\Delta \phi$に等しい仕事が必要である．
        
            \myspaceb%
            $\displaystyle%
                W_{\mathrm{PF}} = 1.6 \times 10^{-19}\:\mathrm{C} \times (14-6)\:\mathrm{V}%
                \fallingdotseq \uwave{1.3 \times 10^{-18}\:\mathrm{J}} \;,
            $

            \myspaceb%
            $\displaystyle%
                W_{\mathrm{PG}} = 1.6 \times 10^{-19}\:\mathrm{C} \times (6-6)\:\mathrm{V}%
                = \uwave{0\:\mathrm{J}} \;,
            $

            \myspaceb%
            $\displaystyle%
                W_{\mathrm{PH}} = 1.6 \times 10^{-19}\:\mathrm{C} \times (4-6)\:\mathrm{V}%
                = \uwave{-3.2 \times 10^{-19}\:\mathrm{J}} \;.
            $
        \item%
            運ぶ際の電位の変化を$\Delta \phi$とすれば，電子の位置エネルギー変化%
            $(-e) \Delta \phi$に等しい仕事が必要である．
        
            \myspaceb%
            $\displaystyle%
                \begin{myaligna}
                    W^*_{\mathrm{PF}}%
                    & = -1.6 \times 10^{-19}\:\mathrm{C} \times (14-6)\:\mathrm{V} \\
                    & \fallingdotseq \uwave{-1.3 \times 10^{-18}\:\mathrm{J}} \;,
                \end{myaligna}
            $

            \myspaceb%
            $\displaystyle%
                W^*_{\mathrm{PG}} = -1.6 \times 10^{-19}\:\mathrm{C} \times (6-6)\:\mathrm{V}
                = \uwave{0\:\mathrm{J}} \;,
            $

            \myspaceb%
            $\displaystyle%
                W^*_{\mathrm{PH}} = -1.6 \times 10^{-19}\:\mathrm{C} \times (4-6)\:\mathrm{V}
                = \uwave{3.2 \times 10^{-19}\:\mathrm{J}} \;.
            $
        \item%
            電気力線は等電位線に直交し，高電位から低電位へ向かう向きである．\\
            \begin{minipage}{\linewidth}
                \centering
                    \includegraphics%
                    {\subfix{fig_uem_phy_std_em_ef_04/fig_uem_phy_std_em_ef_04_03_a.pdf}}
            \end{minipage}
        \vspace{1pt}
        \item%
            電場は電位の空間勾配に等しいので，
            
            \myspaceb%
            $\displaystyle%
                E_\mathrm{P} \fallingdotseq \frac{2\:\mathrm{V}}{0.53 \:\mathrm{m}}
                = 3.77  \cdots \:\mathrm{V/m}
                \fallingdotseq \uwave{3.8 \:\mathrm{V/m}} \;.
            $
    \end{enumerate}
\end{document}

